\documentclass[11pt,a4paper]{ctexart}
\usepackage[margin=2.5cm]{geometry}
\usepackage{amsmath,amssymb,bm}
\usepackage{tikz}
\setlength{\parskip}{.4em}
\title{分析力学第二次作业答案}
\author{张溢致\quad 致理数41\quad 2024011936}
\date{}
\begin{document}
\maketitle
\noindent\textbf{1.} 利用拉格朗日方程做1.5题：质量为 $M$，倾角为 $\theta$ 的光滑斜面的底面放在光滑的水平面上，有一质量为 $m$ 的小块从顶端沿此斜面滑下。给出小块相对斜面的加速度。

\textbf{解：} 设斜面向右下倾斜，取斜面水平位移 $X$ 和小块沿斜面向下的位移 $s$ 为广义坐标。小块的坐标可以写成
\[
x=X+s\cos\theta,\qquad y=h-s\sin\theta.
\]
因此
\[
T=\frac12M\dot X^2+\frac12m\bigl[(\dot X+\dot s\cos\theta)^2+\dot s^2\sin^2\theta\bigr],
\qquad V=mg(h-s\sin\theta).
\]
略去常量 $-mgh$，得到
\[
L=\frac12(M+m)\dot X^2+m\dot X\dot s\cos\theta+\frac12m\dot s^2+mgs\sin\theta.
\]
$X$ 是循环坐标，因为 $L$ 中包含 $\dot X$，却不包含 $X$ 本身。对 $X$ 使用拉格朗日方程，有
\[
\frac{\mathrm d}{\mathrm dt}\frac{\partial L}{\partial\dot X}
-\frac{\partial L}{\partial X}=0.
\]
其中
\[
\frac{\partial L}{\partial\dot X}=(M+m)\dot X+m\dot s\cos\theta,
\qquad\frac{\partial L}{\partial X}=0.
\]
由于质量和倾角均为常量，再对时间求导，得到
\[
(M+m)\ddot X+m\ddot s\cos\theta=0,
\]
对 $s$ 使用拉格朗日方程，有
\[
\frac{\mathrm d}{\mathrm dt}\frac{\partial L}{\partial\dot s}
-\frac{\partial L}{\partial s}=0.
\]
分别求偏导，得到
\[
\frac{\partial L}{\partial\dot s}=m\dot X\cos\theta+m\dot s,
\qquad\frac{\partial L}{\partial s}=mg\sin\theta.
\]
因此
\[
m\ddot X\cos\theta+m\ddot s-mg\sin\theta=0,
\]
除以 $m$ 后为
\[
\ddot s+\ddot X\cos\theta=g\sin\theta.
\]
由关于 $X$ 的方程，得
\[
\ddot X=-\frac{m\cos\theta}{M+m}\ddot s.
\]
将其代入关于 $s$ 的方程，消去 $\ddot X$：
\[
\ddot s-\frac{m\cos^2\theta}{M+m}\ddot s=g\sin\theta.
\]
乘以 $M+m$ 并合并，得到
\[
\bigl[M+m-m\cos^2\theta\bigr]\ddot s
=(M+m)g\sin\theta.
\]
由 $1-\cos^2\theta=\sin^2\theta$，有
\[
\bigl[M+m\sin^2\theta\bigr]\ddot s=(M+m)g\sin\theta.
\]
所以，小块相对斜面的加速度为
\[
\ddot s=\frac{(M+m)g\sin\theta}{M+m\sin^2\theta}.
\]
方向沿斜面向下。斜面本身的加速度为
\[
\ddot X=-\frac{mg\sin\theta\cos\theta}{M+m\sin^2\theta}.
\]
当 $M\to\infty$ 时，$\ddot s\to g\sin\theta$，回到固定斜面的结果。

\noindent\textbf{2.} 两质点质量分别为 $m_1$ 和 $m_2$，坐标为 $\bm r_1$ 和 $\bm r_2$，它们之间的相互作用势能为 $V(x_1-x_2,y_1-y_2,z_1-z_2)$。

(1) 以两质点的质心坐标和相对坐标作为广义坐标，给出拉格朗日量，并找到循环坐标和循环积分；

(2) 给出罗斯函数，并导出运动微分方程。

\textbf{解：} 记总质量和约化质量为
\[
M=m_1+m_2,\qquad \mu=\frac{m_1m_2}{M},
\]
选择
\[
\bm R=\frac{m_1\bm r_1+m_2\bm r_2}{M},\qquad \bm r=\bm r_1-\bm r_2.
\]
反解为
\[
\bm r_1=\bm R+\frac{m_2}{M}\bm r,\qquad
\bm r_2=\bm R-\frac{m_1}{M}\bm r.
\]
代入动能后交叉项抵消，所以
\[
L=\frac M2\dot{\bm R}^{2}+\frac\mu2\dot{\bm r}^{2}-V(\bm r).
\]
其中矢量速度的平方表示自身的点积，例如 $\dot{\bm R}^{2}=\dot{\bm R}\cdot\dot{\bm R}$。
质心的三个分量均为循环坐标，相应的循环积分是总动量
\[
{\bm P=\frac{\partial L}{\partial\dot{\bm R}}=M\dot{\bm R}=\text{常矢量}}.
\]
对一般的 $V(\bm r)$，不能再断言相对坐标也是循环坐标。

采用以下罗斯函数定义：
\[
\mathcal R=L-\bm P\cdot\dot{\bm R}.
\]
将 $\dot{\bm R}=\bm P/M$ 代入此定义，而非仅代入 $L$，得到
\[
\mathcal R=\frac{\bm P^2}{2M}+\frac\mu2\dot{\bm r}^{2}-V(\bm r)-\frac{\bm P^2}{M}
=\frac\mu2\dot{\bm r}^{2}-V(\bm r)-\frac{\bm P^2}{2M}.
\]
这里 $\bm P$ 作为固定参数。对非循环坐标应用拉格朗日型方程，有
\[
\frac{\mathrm d}{\mathrm dt}\frac{\partial\mathcal R}{\partial\dot r_i}-\frac{\partial\mathcal R}{\partial r_i}=0,
\qquad{\mu\ddot{\bm r}=-\nabla_{\bm r}V}.
\]
质心运动由循环积分还原为
\[
{\bm R(t)=\bm R(0)+\frac{\bm P}{M}t}.
\]
若将罗斯函数定义为 $\bm P\cdot\dot{\bm R}-L$，则整个函数及相应方程同时反号，运动结果相同。

\noindent\textbf{3.} 有一质点质量为 $m$，在中心力场 $V(r)$ 中运动，柱坐标为 $(\rho,\varphi,z)$，其拉格朗日量为
\[
L=\frac m2(\dot\rho^2+\rho^2\dot\varphi^2+\dot z^2)-V(\sqrt{\rho^2+z^2}).
\]
另有一个转动坐标系，绕 $z$ 轴以角速度 $\omega$ 转动。若 $t=0$ 时两个坐标系重合，转动坐标系质点坐标为 $(\rho',\varphi',z')$。

(1) 给出转动坐标系坐标表示的拉格朗日量和循环积分；(2) 给出拉格朗日方程；(3) 考查此时的“虚拟力”的形式。

\textbf{解：} 取 $\omega$ 为常量，则
\[
\rho=\rho',\qquad z=z',\qquad\varphi=\varphi'+\omega t.
\]
记 $r=\sqrt{\rho'^2+z'^2}$，有
\[
{L'=\frac m2\bigl[\dot\rho'^2+\rho'^2(\dot\varphi'+\omega)^2+\dot z'^2\bigr]-V(r)}.
\]
$\varphi'$ 是循环坐标，因此
\[
{p_{\varphi'}=m\rho'^2(\dot\varphi'+\omega)=J_z=\text{常量}}.
\]
由于 $L'$ 不显含时间，还存在广义能量积分
\[
\mathcal E'=\frac m2(\dot\rho'^2+\rho'^2\dot\varphi'^2+\dot z'^2)+V(r)-\frac m2\omega^2\rho'^2
=E-\omega J_z.
\]
其中 $E$ 为惯性系中的机械能。

分别对三个坐标列方程，得到
\[
{\begin{aligned}
m\ddot\rho'-m\rho'(\dot\varphi'+\omega)^2&=-V'(r)\frac{\rho'}r,\\
\frac{\mathrm d}{\mathrm dt}\bigl[m\rho'^2(\dot\varphi'+\omega)\bigr]&=0,\\
m\ddot z'&=-V'(r)\frac{z'}r.
\end{aligned}}
\]
将前两式整理成转动系柱坐标加速度的形式：
\[
\begin{aligned}
m(\ddot\rho'-\rho'\dot\varphi'^2)&=-V'(r)\frac{\rho'}r+2m\omega\rho'\dot\varphi'+m\omega^2\rho',\\
m(\rho'\ddot\varphi'+2\dot\rho'\dot\varphi')&=-2m\omega\dot\rho'.
\end{aligned}
\]
令 $\bm\Omega=\omega\bm e_z$，上述额外项正是
\[
{\bm F_{\rm C}=-2m\bm\Omega\times\bm v',\qquad
\bm F_{\rm cf}=-m\bm\Omega\times(\bm\Omega\times\bm r')}.
\]
前者为科里奥利力，后者为离心力，径向与方位向分量分别为
\[
F_\rho^{\rm fict}=2m\omega\rho'\dot\varphi'+m\omega^2\rho',\qquad
F_\varphi^{\rm fict}=-2m\omega\dot\rho',\qquad F_z^{\rm fict}=0.
\]
这里转动角速度恒定，故没有欧拉惯性力。

\noindent\textbf{4.} 利用泛函变分极值方法证明，在给定球面上，两点之间的最短距离是沿着大圆（半径与球体一致的圆）的一段弧线。

\textbf{解：} 设球半径为 $R$，曲线以弧长 $s$ 参数化，位置矢量为 $\bm r(s)$。在球面约束 $\bm r^2=R^2$ 下，使弧长泛函取驻值。引入乘子，考虑
\[
\mathcal I=\int\left[|\dot{\bm r}|+\frac{\lambda}2(\bm r^2-R^2)\right]\mathrm du.
\]
固定端点的欧拉方程是
\[
\frac{\mathrm d}{\mathrm du}\left(\frac{\dot{\bm r}}{|\dot{\bm r}|}\right)=\lambda\bm r.
\]
在该方程中改用弧长参数，乘子相应重定义，得到
\[
\bm r''(s)=\lambda(s)\bm r(s),\qquad |\bm r'(s)|=1.
\]
由 $\bm r^2=R^2$ 两次求导，有 $\bm r\cdot\bm r'=0$ 和 $\bm r\cdot\bm r''=-1$，故
\[
\lambda=-\frac1{R^2},\qquad \bm r''+\frac1{R^2}\bm r=0.
\]
所以曲线具有形式
\[
\bm r(s)=\bm a\cos(s/R)+\bm b\sin(s/R),
\quad |\bm a|=|\bm b|=R,\quad \bm a\cdot\bm b=0.
\]
曲线位于过球心、由 $\bm a,\bm b$ 张成的平面内，因此必为大圆的一段。

还需说明哪一段给出全局最短距离。设两端点的较小球心角为 $\gamma\in[0,\pi]$，以起点方向为球坐标极轴。对任意连接曲线，球面线元为
\[
\mathrm ds=R\sqrt{\mathrm d\vartheta^2+\sin^2\vartheta\,\mathrm d\phi^2}
\ge R|\mathrm d\vartheta|.
\]
因此总长度至少为 $R\gamma$。连接端点的大圆短弧恰好达到此下界，故
\[
{\ell_{\min}=R\gamma}.
\]
两端点非对径时短弧唯一；对径时，每条连接二者的大圆半圆弧都最短。

\noindent\textbf{5.} 地球表面大气由于重力导致密度不均，也因此造成光的折射率随高度变化。若气体分子数密度可以用 $n(y)$ 来表示，$y$ 是距地面的垂直高度，则 $n(y)=n_0e^{-\alpha y}$，$n_0$ 是地面上气体分子数密度，而折射率为 $\sqrt{1+\beta n(y)}$。为简化问题，可以假设 $\alpha$、$\beta$ 都是常量。由于大气厚度相比地球半径小很多，可以忽略地面弧度，只考虑高度变化。当太阳光线以某个入射角 $\theta$ 进入大气层后，光线通过大气到达地面的轨迹可以通过光程极小来确定，也就是使得积分值
\[
\int_{\infty}^{0}\sqrt{1+\beta n(y)}\,\mathrm ds
\]
极小，这里 $\mathrm ds$ 是光线的长度。试利用泛函变分极值方法给出光线轨迹方程。

\textbf{解：} 将分子数密度与折射率区别记为
\[
a=\beta n_0,\qquad N(y)=\sqrt{1+ae^{-\alpha y}}.
\]
按题意取 $\alpha>0$、$a\ge0$，入射角从竖直法线量起。先在有限高度 $H$ 与地面之间固定端点，再取 $H\to\infty$。这避免把无限光程直接当作有限泛函。

在光线所在的竖直平面内，将轨迹写成 $x=x(y)$，有限段的正光程为
\[
\mathcal J[x]=\int_0^H N(y)\sqrt{1+x'^2}\,\mathrm dy.
\]
被积函数不含 $x$，因此欧拉方程给出
\[
\frac{N(y)x'}{\sqrt{1+x'^2}}=C.
\]
选择水平轴使反向追迹（从地面向上）时 $x'\ge0$。无穷高处 $N\to1$、$x'\to\tan\theta$，从而 $C=\sin\theta$。于是
\[
{\frac{\mathrm dx}{\mathrm dy}=\frac{\sin\theta}{\sqrt{\cos^2\theta+ae^{-\alpha y}}}}.
\]
若地面交点为 $(x_g,0)$，轨迹可直接写成
\[
{x(y)=x_g+\sin\theta\int_0^y\frac{\mathrm d\eta}{\sqrt{\cos^2\theta+ae^{-\alpha\eta}}}}.
\]
若采用相反的水平轴，右端积分前改为负号即可。

当 $a>0$、$0<\theta<\pi/2$ 时，记 $c=\cos\theta$，令 $u=\sqrt{c^2+ae^{-\alpha y}}$，则
\[
\frac{\mathrm dy}{u}=-\frac2\alpha\frac{\mathrm du}{u^2-c^2}.
\]
积分得到显式轨迹
\[
{x(y)=x_g+\frac{2\tan\theta}{\alpha}
\left[\operatorname{artanh}\frac{c}{\sqrt{c^2+ae^{-\alpha y}}}
-\operatorname{artanh}\frac{c}{\sqrt{c^2+a}}\right]}.
\]
沿途与竖直方向的夹角 $i(y)$ 满足
\[
{N(y)\sin i(y)=\sin\theta},\qquad
\sin i(0)=\frac{\sin\theta}{\sqrt{1+a}}.
\]
因折射率向地面增大，光线逐渐向竖直法线偏折。$\theta=0$ 时为竖直直线，$a=0$ 时为 $x=x_g+y\tan\theta$。积分形式还覆盖掠入射的极限。

\noindent\textbf{6.} 质量为 $m$ 的粒子的势能，在半径为 $R$ 的球形区域内为 $-|U_0|$，球外为零。该粒子以初速度 $v_0$ 从球面某处射入球内，之后离开球。根据哈密顿原理给出运动轨迹。

\textbf{解：} 记 $U=|U_0|>0$。球内外势能分别为常量，故粒子在各区域都作匀速直线运动。能量守恒给出
\[
E=\frac12mv_0^2,\qquad
v_1=\sqrt{v_0^2+\frac{2U}{m}},\qquad
q=\frac{v_0}{v_1}<1.
\]
关键在于球面处的连接条件。由哈密顿原理，对分段轨迹的穿越时刻及穿越点作变分，作用量边界项为
\[
\delta S=(\bm p_- -\bm p_+)\cdot\delta\bm r
-(H_- -H_+)\delta t.
\]
其中 $\delta\bm r$ 沿球面切向任意，$\delta t$ 也任意。因此界面两侧的切向动量相等、能量相等：
\[
\bm p_{-,\parallel}=\bm p_{+,\parallel},\qquad H_-=H_+.
\]
设入射角 $i$ 和球内折射角 $r$ 都从入射点指向球心的法线量起，则
\[
{v_0\sin i=v_1\sin r},\qquad
{\sin r=q\sin i}.
\]
因此进入球内时向法线偏折。这里粒子界面的折射规律由动量条件给出。

设入射直线对球心的瞄准距离（冲量参数）为 $b\in[0,R)$，则
\[
\sin i=\frac bR,\qquad
\sin r=\frac{qb}{R}.
\]
角动量守恒使运动始终位于入射方向与球心共同确定的平面内。球内直线到球心的距离为 $qb$，弦长为
\[
\ell=2R\cos r=2\sqrt{R^2-q^2b^2}.
\]
弦与两个端点的半径夹角相等，所以球内到达出射点的入射角仍为 $r$。出射时
\[
v_1\sin r=v_0\sin i_{\rm out},
\qquad i_{\rm out}=i.
\]
粒子恢复速率 $v_0$，背离出射法线偏折。由于该轨迹的 $\sin r=qb/R<q$，它能正常出射。

若取入射方向沿 $+x$，入射直线为 $y=b$，令进入球面时 $t=0$，定义
\[
i=\arcsin(b/R),\quad r=\arcsin(qb/R),\quad\delta=i-r,
\]
则入射点和球内单位方向为
\[
\bm A=(-R\cos i,R\sin i),\qquad
\bm u=(\cos\delta,-\sin\delta).
\]
出射点与出射时刻为
\[
\bm B=\bm A+2R\cos r\,\bm u,\qquad
\tau=\frac{2R\cos r}{v_1}.
\]
完整轨迹是三段直线：
\[
{\bm x(t)=\begin{cases}
\bm A+v_0t(1,0),&t\le0,\\
\bm A+v_1t\bm u,&0\le t\le\tau,\\
\bm B+v_0(t-\tau)(\cos2\delta,-\sin2\delta),&t\ge\tau.
\end{cases}}
\]
入射与出射方向的偏转角大小为
\[
{\chi=2(i-r)=2\left[\arcsin\frac bR-
\arcsin\left(\frac{v_0b}{Rv_1}\right)\right]}.
\]
方向朝球心一侧。$b=0$ 时沿直径直进直出；$U\to0$ 时 $q\to1$、$\chi\to0$，回到一条直线。
\end{document}
